\documentclass[11pt]{article}

\usepackage[margin=1in]{geometry}
\usepackage{amsmath,amssymb,amsthm,mathtools}
\usepackage{booktabs,tabularx,array}
\usepackage{microtype}
\usepackage{enumitem}
\usepackage{xcolor}
\usepackage{xurl}
\usepackage[hidelinks]{hyperref}

% Drafted 2026-07-11 from Sources/Furey_Baez_DVT_Formalization_Paper_Outline.md
% (2026-05-31), updated against the live trusted tree: every theorem cited
% below was checked against its module docstring and, for the flagships,
% against the build-enforced axiom guards on the day of writing.  Venue
% target per the outline: Annals of Formalized Mathematics (primary),
% ITP-style cut in reserve.  Author line and release gates are user
% decisions.

\newcommand{\ii}{\mathrm{i}}
\newcommand{\id}{\mathbf{1}}
\newcommand{\C}{\mathbb{C}}
\newcommand{\R}{\mathbb{R}}
\newcommand{\OO}{\mathbb{O}}
\newcommand{\su}{\mathfrak{su}}
\definecolor{kernelcolor}{RGB}{0,92,74}
\newcommand{\Kernel}{\textcolor{kernelcolor}{\textsc{Kernel}}}
\newcommand{\leanname}[1]{\texttt{#1}}

\newtheorem{theorem}{Theorem}
\newtheorem{proposition}{Proposition}
\newtheorem{remark}{Remark}

\title{Verified Octonionic Algebra for Standard Model Gauge Structure:\\
a Lean~4 formalization of the Furey and
Baez--Dubois-Violette--Todorov programs}
\author{Mark Schwab}
\date{July 11, 2026 --- draft v1}

\begin{document}
\maketitle

\begin{abstract}
Several proposed links between the octonions and the Standard Model rely on
closely related algebraic structures: a chosen complex subalgebra of
$\OO$, Clifford algebras generated by complex-octonionic ladder operators,
the exceptional Jordan algebra $h_3(\OO)$, and the compact gauge group
$S(U(2)\times U(3)) \cong (SU(3)\times SU(2)\times U(1))/\mathbb{Z}_6$.  We
report a Lean~4 formalization, over Mathlib, of a common algebraic core
underlying the programs of Furey and of
Baez--Dubois-Violette--Todorov, with every theorem kernel-checked at the
standard axiom footprint and with build-enforced axiom guards on the
flagships.  On the Baez side we prove that the group of octonion algebra
automorphisms fixing the distinguished imaginary unit is multiplicatively
equivalent to $SU(3)$ --- with the target literally Mathlib's
$\leanname{Matrix.specialUnitaryGroup}$ --- together with
Pati--Salam/Standard-Model/$SU(5)$ block-embedding equivalences and an
exact six-element kernel and unit-level quotient-group package for the
$\mathbb{Z}_6$ covering of $S(U(2)\times U(3))$.  On the Furey side we
formalize the ladder operators in an explicit convention bridge, their
$\mathrm{Cl}(6)$ anticommutation relations, the primitive idempotent and
its minimal left ideal, the octonionic charge operator and its eigenvalue
table, the operator-level Gell-Mann--Nishijima identity, an electroweak
package with $\su(2)$ raising and lowering operators, anomaly sums tied to
the charge-operator eigenvalues, and a one-generation realization package
whose remaining gap --- the conventional right-handed singlet completion
--- is recorded as a machine-readable claim boundary.  On the Jordan side
we prove the Jordan identity for a coordinate model of $h_3(\OO)$, its
trace-form splitting over the chosen $h_3(\C)$, and a full algebraic
characterization of the Dubois-Violette--Todorov two-sided
$(SU(3)\times SU(3)^{\mathrm{op}})/\mathbb{Z}_3$ stabilizer action on the
complement.  A reusable family-symmetry interface transports charge
tables, eigenvectors, and anomaly sums along any family action that
commutes with the charge data, instantiated for Furey--Hughes triality
role permutations.  The contribution is not a physical derivation of the
Standard Model: it is a verified, convention-explicit audit trail through
constructions that are usually compared only informally, with the
unproved steps stated as visibly as the proved ones.
\end{abstract}

\section{Introduction}

The octonions keep reappearing in proposals about Standard Model
structure: as the source of a distinguished $SU(3)$
\cite{Baez2021,BaezOctonions}, as a generator of one fermion generation
through ladder operators acting on a minimal left ideal
\cite{Furey2016,Furey2018}, as the coordinate ring of the exceptional
Jordan algebra whose substructures select
$S(U(2)\times U(3))$ \cite{DVT2018}, and as the organizing algebra behind
recent triality-based accounts of the three generations
\cite{FureyHughes2025}.  These proposals share more than a theme.  They
share concrete algebraic objects --- a chosen complex line in $\OO$, the
induced splitting $\OO \cong \C \oplus \C^3$, operators built from left
multiplication, and the same compact gauge-group target --- and they are
usually compared informally, across papers with incompatible basis
conventions, implicit parenthesizations, and different standards for what
counts as ``deriving'' Standard Model data.

This paper reports a formalization of that shared core in Lean~4 over
Mathlib.  The novelty is not a new physical derivation.  It is a verified
mathematical audit trail: a single repository in which the Furey ladder
construction, the Baez stabilizer route, the
Dubois-Violette--Todorov (DVT) Jordan route, and the conventional
$\mathbb{Z}_6$ gauge quotient are all built on one explicitly fixed
octonion convention, with kernel-checked theorems where the informal
literature has calculations, and machine-readable claim boundaries where
the informal literature has hopes.

\paragraph{Contributions.}
\begin{enumerate}[leftmargin=2em]
\item A convention-explicit octonion model with a verified translation
  bridge to the Baez/Furey basis conventions, isolating the single sign
  flip that separates them (Section~\ref{sec:conventions}).
\item The Baez stabilizer theorem at the algebraic level: the group of
  octonion algebra automorphisms fixing the distinguished imaginary unit
  $e_{111}$ is multiplicatively equivalent to $SU(3)$, where the target
  is definitionally Mathlib's special unitary group
  (Section~\ref{sec:baez}).
\item A trusted formalization of Furey's complex-octonion ladder
  operators, their $\mathrm{Cl}(6)$ relations, the primitive idempotent
  and minimal left ideal, the charge operator, the operator-level
  Gell-Mann--Nishijima identity, an electroweak $\su(2)$ package, anomaly
  sums derived from the charge operator, and a one-generation realization
  package with an explicit right-handed-sector boundary
  (Section~\ref{sec:furey}).
\item An exact algebraic $\mathbb{Z}_6$ package for the Standard Model
  gauge group: the covering homomorphism onto the block group, its
  six-element kernel with a product-level identity-fiber iff, and a
  unit-level quotient-group equivalence (Section~\ref{sec:gauge}).
\item A coordinate formalization of $h_2(\OO)$ and $h_3(\OO)$ including
  the Jordan identity, trace-form splitting over $h_3(\C)$, and a full
  algebraic characterization of the DVT two-sided
  $(SU(3)\times SU(3)^{\mathrm{op}})/\mathbb{Z}_3$ stabilizer action on
  the complement (Section~\ref{sec:jordan}).
\item A reusable family-symmetry naturality interface --- transport of
  charge tables, eigenvectors, and anomaly sums along commuting family
  actions --- instantiated for Furey--Hughes triality role permutations
  (Section~\ref{sec:family}).
\item A documented claim-boundary methodology for formalizing
  convention-heavy, physically speculative mathematics without
  overclaiming (Sections~\ref{sec:conventions} and
  \ref{sec:boundaries}).
\end{enumerate}

\paragraph{Claim boundary, stated first.}  Nothing here derives the
Standard Model.  The formalization verifies finite algebraic and
coordinate-Jordan statements that the octonionic programs rely on, and it
verifies them in one shared convention.  It does not prove
$\mathrm{Aut}(h_3(\OO)) \cong F_4$, does not prove the full DVT compact
stabilizer-intersection theorem, does not prove topological or smooth
statements about any of the groups involved, does not derive the
right-handed fermion sector from the ladder algebra, and does not touch
dynamics, Lagrangians, masses, or mixing.  Section~\ref{sec:boundaries}
lists the non-results as explicitly as the results.

\section{Related work}
\label{sec:related}

\paragraph{Octonions and the Standard Model.}  The mathematical
background is Baez's survey of the octonions \cite{BaezOctonions}.  The
specific $SU(3)$ stabilizer route formalized here follows Baez's 2021
account \cite{Baez2021}; the grand-unified block structure follows
Baez--Huerta \cite{BaezHuerta2009}.  Furey's program
\cite{Furey2015,Furey2016,Furey2018} builds one generation of fermions
from $\C\otimes\OO$ acting on itself; the recent Furey--Hughes work
\cite{FureyHughes2025} adds a triality-based account of three
generations.  Dubois-Violette and Todorov \cite{DVT2018} obtain the
Standard Model gauge group from stabilizers inside the automorphism group
of the exceptional Jordan algebra $h_3(\OO)$; Boyle \cite{Boyle2020}
and Krasnov \cite{Krasnov2019,Krasnov2025} develop closely related
complex-structure and triality perspectives, and Gresnigt
\cite{Gresnigt2026} pursues an $S_3$-symmetric three-generation model in
$\mathrm{Cl}(10)$.  Our formalization treats these as sources of exact
finite statements to verify, not as physical claims to adjudicate.

\paragraph{Formalized physics.}  The closest neighbor is
PhysLean/HepLean \cite{ToobySmith2025,ToobySmith2024}, which digitalizes
high-energy-physics definitions and calculations in Lean~4 (CKM matrices,
local anomaly cancellation conditions, Higgs-sector material, index
notation).  The present work is complementary: it formalizes octonionic
and exceptional-Jordan \emph{algebraic routes to the gauge-group
scaffold} rather than phenomenology-facing computations.  Where the two
touch --- anomaly arithmetic --- our anomaly-freedom predicates are
self-contained finite statements, and interfacing them with
PhysLean-style anomaly packages is natural future work.  We consulted
PhysLean throughout as a version-pinned reference library and cross-check
target rather than as an import.

\section{Conventions, architecture, and trust}
\label{sec:conventions}

All results live in a single Lean~4 repository pinned to
\leanname{leanprover/lean4:v4.28.0} over Mathlib, organized into a trusted
layer (\leanname{PhysicsSM/Algebra}, \leanname{PhysicsSM/Gauge},
\leanname{PhysicsSM/StandardModel}) and an explicitly separated draft
layer not cited in this paper.  Trusted modules build with no
\leanname{sorry}-like placeholders and no new axioms.  The flagship
theorems carry \emph{build-enforced axiom guards}
(\leanname{Furey.AxiomGuard}): \leanname{\#guard\_msgs} blocks around
\leanname{\#print axioms} that fail the build if any placeholder, native
token, or new axiom ever leaks into their transitive dependencies.
Every other cited theorem was checked, at the time of writing, to carry
the standard Mathlib logical footprint
(\leanname{propext}, \leanname{Classical.choice}, \leanname{Quot.sound})
and nothing else --- in particular, none of the results cited in this
paper use compiler-trusting evaluation (\leanname{native\_decide}) ---
but only the flagships have that footprint enforced on every build; the
remaining checks are reproducible one-line \leanname{\#print axioms}
queries.

\paragraph{The octonion convention.}  The project fixes an explicit basis
$e_{000}$ (the unit), $e_{001}, e_{010}, e_{011}, e_{100}, e_{101},
e_{110}, e_{111}$, with product index given by bitwise XOR of the labels
and signs fixed by a Fano-plane orientation
(\leanname{Algebra.Octonion.Basic}).  This convention is \emph{not}
Baez's or Furey's.  The translation is itself a verified object
(\leanname{Algebra.Octonion.ConventionBridge}): a relabeling permutation
together with exactly one sign flip,
\[
e_1 \mapsto e_{001},\quad e_2 \mapsto e_{010},\quad e_3 \mapsto e_{110},
\quad e_4 \mapsto e_{011},\quad
\boxed{e_5 \mapsto -e_{100}},\quad e_6 \mapsto e_{101},\quad
e_7 \mapsto e_{111},
\]
under which all $21$ oriented Fano products and their $21$ reversals are
verified to agree with an independent external oracle.  Every formula
imported from the literature passes through this bridge; no Baez or Furey
structure constant enters the repository verbatim.  Nonassociativity is
handled by discipline rather than convention: octonion products are
explicitly parenthesized terms, and operator statements are phrased in
terms of left-multiplication maps, which \emph{do} compose associatively.

\begin{center}
\begin{tabular}{lll}
\toprule
Issue & Informal-literature risk & Formalization choice \\
\midrule
Fano orientation & sign drift between papers & fixed XOR convention + bridge \\
Nonassociativity & omitted parentheses & explicit Lean terms; operator algebra \\
Gauge group & product vs.\ quotient & $S(U(2)\times U(3))$ and exact
$\mathbb{Z}_6$ kernel \\
Interpretation & overclaiming & trusted/draft split; machine-readable
boundaries \\
\bottomrule
\end{tabular}
\end{center}

\paragraph{What the kernel checks.}  Kernel acceptance certifies that a
formal statement follows from its formal hypotheses.  It does not certify
that the statement is the right abstraction of the physics; that judgment
is made in prose here, with the formal statement quoted so the reader can
audit the correspondence.  Proof scripts were partly produced by a proof
agent and were accepted only through the kernel; Section~\ref{sec:eng}
describes the workflow.

\section{The octonionic foundation}
\label{sec:foundation}

The constructions all start from one choice: the imaginary unit
$e_{111}$ and the complex line
$\C \;\cong\; \mathrm{span}_\R\{1, e_{111}\} \subset \OO$.
The trusted foundation layer (\leanname{Algebra.Octonion.ComplexLine},
\leanname{ComplexSplitting}, \leanname{ChosenComplexAlgebra},
\leanname{ComplexTripleModule}) verifies that this two-dimensional real
subspace is closed under multiplication and conjugation, that left
multiplication by $e_{111}$ squares to $-1$ on the complement, and that
the complement carries the structure of a three-dimensional complex
coordinate space $\C^3$ with a Hermitian form
(\leanname{ComplexTripleHermitian}).  This gives the splitting
\[
\OO \;\cong\; \C \,\oplus\, \C^3
\]
on which everything else is built: the Furey ladder operators use the
complexification $\C\otimes\OO$, the Baez route acts on the $\C^3$
complement, and the Jordan route embeds $h_3(\C)$ inside $h_3(\OO)$ by
applying the same splitting entrywise.

\section{The Baez route: automorphisms fixing $e_{111}$ are exactly $SU(3)$}
\label{sec:baez}

Baez's observation \cite{Baez2021} is that the subgroup of $G_2 =
\mathrm{Aut}(\OO)$ stabilizing a unit imaginary octonion is $SU(3)$, and
that this is the natural octonionic home of the color group.  The
formalization proves the algebraic content of this statement exactly,
in the following form.

\begin{theorem}[Stabilizer equivalence, algebraic level]
\label{thm:su3}
Let $\mathrm{Aut}_{e_{111}}(\OO)$ denote the group of multiplicative
linear automorphisms of the project octonion model fixing $1$ and
$e_{111}$.  Restriction to the complement $\C^3$ induces a map that is:
(i) well defined --- every such automorphism preserves the Hermitian form
and acts $\C$-linearly on the complement; (ii) valued in determinant-one
unitaries --- multiplicativity applied to the octonion relation
$e_{001}e_{010} = e_{011}$ pins
$\det = 1$ (\leanname{G2FixingE111DetOne}); (iii) a group homomorphism with trivial kernel; and (iv)
surjective --- every $SU(3)$ matrix extends, uniquely, to an automorphism
fixing the complex line.  Consequently
\[
\mathrm{Aut}_{e_{111}}(\OO) \;\simeq^*\; SU(3),
\]
a multiplicative equivalence onto the submonoid that is proved
\emph{equal} to Mathlib's
\leanname{Matrix.specialUnitaryGroup (Fin 3) C}.  Since a multiplicative
equivalence between groups preserves inverses automatically, this is a
group isomorphism; the packaged Lean object is the multiplicative
equivalence, and the inverse-preservation transport is the standard
one-line consequence, not a separately formalized theorem.
\end{theorem}
\noindent\Kernel{}\quad
(\leanname{G2HermitianPreservation};
\leanname{G2FixingE111DetOne};
\leanname{G2FixingE111Composition};
\leanname{G2FixingE111SU3Equiv.su3MatrixExtension};
\leanname{G2AutomorphismSU3Exactness} ---
\leanname{toSU3Hom} bijective, trivial kernel, surjective, with the
bundled \leanname{G2AutomorphismSU3ExactnessPackage}; and the
guard-pinned identification
\leanname{su3Submonoid\_eq\_specialUnitaryGroup}.)

\begin{remark}[Boundary]
The theorem is a statement about the algebraically defined automorphism
group of the explicit octonion model.  It does not formalize the
topological or smooth Lie group $G_2$, its connectedness, or the
statement ``$\mathrm{Stab}_{G_2}(e_{111}) \cong SU(3)$'' as a theorem
about manifolds.  At the level at which the octonionic Standard Model
literature actually uses the fact --- automorphisms act on
$\C^3$ exactly as $SU(3)$ --- it is closed.
\end{remark}

Building on the same splitting, the block-embedding layer
(\leanname{G2C3GUTPaperPackage}, \leanname{G2C3GUTSU5Bridge},
\leanname{Gauge.GUTSquare}) verifies the grand-unified bookkeeping: for a
Hermitian-form-preserving automorphism $g$ with complement matrix $M$,
the block matrix $\mathrm{diag}(\id_2, M)$ satisfies the Pati--Salam
block predicate unconditionally, satisfies the Standard-Model and
$SU(5)$ block predicates exactly when $\det M = 1$, and the latter two
predicates are equivalent for such maps.  This places the octonionic
$SU(3)$ inside the same block group used by the gauge-quotient package of
the next section, completing the ``GUT square'' of
Baez--Huerta \cite{BaezHuerta2009} at the algebraic level.

\section{The gauge endpoint: $S(U(2)\times U(3))$ and the exact
$\mathbb{Z}_6$ kernel}
\label{sec:gauge}

The compact gauge group of the Standard Model is not the product
$SU(3)\times SU(2)\times U(1)$ but its quotient by a $\mathbb{Z}_6$
subgroup, isomorphic to $S(U(2)\times U(3))$
\cite{BaezHuerta2009}.  The formalization records this endpoint
precisely, as algebra.

\begin{theorem}[Exact $\mathbb{Z}_6$ kernel and unit-level quotient]
\label{thm:z6}
Consider the covering homomorphism
\[
\iota:\; U(1)\times SU(2)\times SU(3) \longrightarrow S(U(2)\times U(3)),
\qquad
(\alpha, g, h) \longmapsto (\alpha^3 g,\; \alpha^{-2} h).
\]
The formalization proves: (i) each of the six explicit kernel elements
$(\alpha, \alpha^{-3}\id_2, \alpha^{2}\id_3)$ with $\alpha^6 = 1$ maps to
the identity, and the kernel-element type has cardinality exactly $6$;
(ii) at the product level, where the $SU(2)$ and $SU(3)$ factors carry
individual determinant constraints, the identity fiber is
\emph{precisely} these six elements (an iff); and (iii) the quotient of
the covering triples by the kernel is multiplicatively equivalent to the
block-units subgroup,
$\leanname{SMCoveringQuotient} \simeq^* \leanname{SMBlockUnitsSubgroup}$,
with a bundled unit-level quotient-group package.
\end{theorem}
\noindent\Kernel{}\quad
(\leanname{Gauge.StandardModelUnitZ6ExactKernelPackage};
\leanname{StandardModelProductCoveringTrueZ6Kernel};
\leanname{StandardModelZ6FiniteKernel};
\leanname{StandardModelUnitZ6QuotientGroup} and
\leanname{StandardModelUnitZ6QuotientGroupPackage};
supporting block maps in \leanname{StandardModelBlockHom},
\leanname{BlockEmbeddings}.)

\begin{remark}[Boundary]
These are algebraic (group-theoretic) statements about explicit matrix
groups and their quotients.  No topological quotient, covering-space, or
compact-Lie-group statement is claimed.  The distinction matters for
honesty, not for the way the octonionic programs use the fact: their
bookkeeping is exactly the finite kernel arithmetic verified here.
\end{remark}

\section{The Furey route: ladder operators, the minimal ideal, and one
generation}
\label{sec:furey}

Furey's construction \cite{Furey2016,Furey2018} works in the
complexification $\C\otimes\OO$ and builds one generation of fermions
from the action of left-multiplication operators on a minimal left
ideal.  Because $\OO$ is nonassociative, the formalization is careful
about what acts on what: states are elements of $\C\otimes\OO$; operators
are $\C$-linear maps generated by left multiplications; and ``minimal
left ideal'' always means a minimal left module for that associative
operator algebra, not an ideal of the nonassociative ring itself.

\subsection{Ladder operators and the $\mathrm{Cl}(6)$ relations}

Furey's ladder operators translate through the convention bridge to
\[
\alpha_1 = \tfrac{1}{2}(+e_{100} + \ii e_{011}),\qquad
\alpha_2 = \tfrac{1}{2}(-e_{110} + \ii e_{001}),\qquad
\alpha_3 = \tfrac{1}{2}(-e_{101} + \ii e_{010}),
\]
the sign in $\alpha_1$ being forced by the single sign flip in the
bridge.  The formalization proves the canonical anticommutation
relations for the associated left-multiplication operators: each
$\alpha_k$ is nilpotent, $\{\alpha_i, \alpha_j\} = 0$, and
$\{\alpha_i, \alpha_j^\dagger\} = \delta_{ij}$
(\leanname{LadderOperators}; the flagships
\leanname{alpha1\_nilpotent} and \leanname{anticomm\_1\_1dag} are
guard-pinned).  These are exactly the generator relations of
$\mathrm{Cl}(6)$, which is the algebraic engine of the construction.

\subsection{The idempotent, the ideal, and the state dictionary}

The primitive idempotent is $\omega = \tfrac12(1 - \ii e_{111})$
(project convention; \leanname{MinimalLeftIdeal}), and the minimal left
module is generated from it by the ladder operators, an
eight-dimensional complex space per ideal.  There are two conjugate
ideals, and their occupation dictionaries differ; following the primary
source (Furey, Eur.\ Phys.\ J.\ C 78:375 (2018), Eqs.~(22)--(23) and
(26), verified against the stored full text) they are:
\[
S^u:\ \ \nu\ (0,\,\mathbf{1});\quad
\bar D^{r,g,b}\ (+\tfrac13,\,\bar{\mathbf 3});\quad
U^{r,g,b}\ (+\tfrac23,\,\mathbf 3);\quad
E^{+}\ (+1,\,\mathbf 1)
\]
at ladder degrees $0,1,2,3$ with $Q = +\tfrac13 N$, and
\[
S^d:\ \ \bar\nu\ (0,\,\mathbf{1});\quad
D^{r,g,b}\ (-\tfrac13,\,\mathbf 3);\quad
\bar U^{r,g,b}\ (-\tfrac23,\,\bar{\mathbf 3});\quad
E^{-}\ (-1,\,\mathbf 1)
\]
with $Q = -\tfrac13 N$; the $SU(3)$ content is
$S^u \sim \mathbf 1 \oplus \bar{\mathbf 3} \oplus \mathbf 3 \oplus
\mathbf 1$ and
$S^d \sim \mathbf 1 \oplus \mathbf 3 \oplus \bar{\mathbf 3} \oplus
\mathbf 1$ (Furey Eq.~(26)).  The formalized sector
(\leanname{AnomalyBridge}: \leanname{Q\_op\_omega\_bar},
\leanname{Q\_op\_vbar1}--\leanname{vbar6}, \leanname{Q\_op\_nu\_bar})
carries exactly the $S^d$ pattern: kernel-checked eigenvalues $0$,
$-\tfrac13$ (degree one), $-\tfrac23$ (degree two), $-1$ (degree
three).  The particle dictionary is recorded as a labeling, not a
derivation, and four boundaries are mandatory when reading it.  First,
the two ideals are exchanged by $\ii \mapsto -\ii$; Furey identifies
this map with particle--antiparticle conjugation \emph{within the
model}, and that identification is an interpretive postulate of the
construction, not a theorem --- representation conjugacy and physical
antiparticle conjugation must not be conflated.  Second, the primitive
idempotent that plays the role of a vacuum is an algebraic idempotent,
not the QFT vacuum state.  Third, occupation degree is an algebraic
grading; it carries no claim that leptons are composites of quarks or
similar.  Fourth, nothing at this level makes the empty and full states
a weak doublet; weak structure enters separately
(Section~\ref{sec:furey} electroweak layer), and there it is a
consistent construction on the states, not a derivation from the color
ladder algebra.

\subsection{Charge, electroweak structure, and the operator-level
Gell-Mann--Nishijima identity}

The charge operator $Q_{\mathrm{op}} = \tfrac13(N_1 + N_2 + N_3)$ built
from the ladder number operators acts diagonally on the state basis with
the physical electric charges as eigenvalues
(\leanname{QuantumNumbers}, \leanname{OperatorRepresentations}).  On top
of this the electroweak layer proves, at the operator level
(\leanname{ElectroweakCompletePackage}, bundling
\leanname{FureyElectroweakPaperPackage}, \leanname{T3OpJbarPackage},
\leanname{FureySU2LDoubletPackage},
\leanname{FureyWeakIsospinLadderPackage}):

\begin{itemize}[leftmargin=2em]
\item the $T_3$ operator and its doublet eigenvalues, and the four
  $SU(2)_L$ doublet pairings on the $\bar J$ sector;
\item the raising/lowering operators $W^\pm$ with the $\su(2)$
  commutation relations, and the derived charge relations
  $[Q, W^+] = W^+$, $[Q, W^-] = -W^-$;
\item the Gell-Mann--Nishijima identity $Q = T_3 + Y/2$ for every basis
  state, with $Q$ the octonionic charge-operator eigenvalue
  (\leanname{HyperchargeBridge}, \leanname{QopElectroweakConsistency},
  \leanname{OperatorElectroweakIdentity}).
\end{itemize}

\begin{remark}[Boundary]
The $W^\pm$ operators are \emph{defined} as explicit basis-state
permutation maps and then proved to satisfy the $\su(2)$ relations; they
are not constructed from the $\alpha_i$ ladder algebra.  Deriving the
electroweak $\su(2)$ from the octonionic operators (rather than
exhibiting it on the states) is an open target, and the package records
this in its docstring.  Likewise the $T_3$ and $Y$ \emph{values} enter
from the conventional table; what is proved is the exact compatibility
of that table with the octonionically derived $Q$.
\end{remark}

\subsection{Anomaly cancellation}

The anomaly layer makes a three-step claim, each step exact:

\begin{enumerate}[leftmargin=2em]
\item \emph{Charge sums from the operator.}  The $U(1)$ linear and cubic
  anomaly sums over the eight $J$-sector states are proved to be the sums
  of the $Q_{\mathrm{op}}$ \emph{eigenvalues} --- each conjunct of the
  theorem is the proved eigenvalue equation for one state, and the final
  conjuncts sum exactly those eigenvalues (to $-4$ and $-2$
  respectively).  This ties the anomaly arithmetic to the octonionic
  operator rather than to hard-coded rational literals, repairing a gap
  found by an internal red-team audit and documented in the module
  (\leanname{AnomalyFromQop.u1Anomaly\_sums\_from\_Qop}).
\item \emph{Doublet-sector cancellation.}  The $\bar J$ left-handed
  doublets (quark doublet $Q_L$, lepton doublet $L_L$) match the
  conventional table in color dimension, weak dimension, and
  hypercharge, and satisfy the $SU(2)^2\,U(1)_Y$ anomaly cancellation
  condition (\leanname{JbarElectroweakAnomaly},
  \leanname{ElectroweakAnomalyBridge}).
\item \emph{Full one-generation anomaly freedom.}  The combined table ---
  Furey doublet sector plus the conventional right-handed singlet
  completion, decomposed sector-by-sector in
  \leanname{FureyAnomalyDecomposition} --- equals the standard
  one-generation table \leanname{standardModelOneGeneration}, which is
  proved locally anomaly free (all six coefficients) and to satisfy
  Witten's global $SU(2)$ condition
  (\leanname{OneGenerationPackage}).
\end{enumerate}

\subsection{The one-generation realization theorem, with its boundary}

The capstone (guard-pinned) is
\leanname{fureyRealizesOneGenerationPackage}
(\leanname{FureyRealizesOneGeneration}): the Furey $\bar J$ doublets
match the Standard Model left-handed sector exactly (color dimension,
weak dimension, hypercharge); every $\bar J$ basis state satisfies
$Q = T_3 + Y/2$ with $Q$ the octonion-derived charge; the doublet
sub-table cancels the $SU(2)^2 U(1)_Y$ anomaly; and the all-left table
matches, once the conventional right-handed singlets are appended, the
anomaly-free \leanname{standardModelOneGeneration}.  The same structure
carries an explicit \leanname{ClaimBoundary} field recording, as a
proposition inside the theorem package itself, that the right-handed
singlet sector $(u_R, d_R, e_R, \nu_R)$ is conventional input, not an
output of the ladder algebra.

We regard this design --- the open gap stated as a machine-readable
component of the flagship theorem --- as one of the paper's
methodological contributions.

\section{The Jordan route: $h_3(\OO)$ and the DVT stabilizer action}
\label{sec:jordan}

Dubois-Violette and Todorov \cite{DVT2018} locate the Standard Model
gauge group inside the automorphism group $F_4$ of the exceptional
Jordan algebra $h_3(\OO)$, as the stabilizer of the splitting induced by
a choice of $\C \subset \OO$.  The formalization builds the coordinate
layer of this story on the same chosen complex line as
Sections~\ref{sec:baez}--\ref{sec:furey}.

\begin{theorem}[Coordinate Jordan layer]
The explicit coordinate model of $h_3(\OO)$ (3$\times$3 Hermitian
octonionic matrices with the symmetrized product
$A \circ B = \tfrac12(AB + BA)$) satisfies the Jordan identity
$(a\circ b)\circ(a\circ a) = a\circ(b\circ(a\circ a))$; the spin-factor
model $h_2(\OO)$ carries its verified product and determinant
vocabulary; and the trace form splits $h_3(\OO)$ orthogonally over the
embedded $h_3(\C)$ with verified complement projections.
\end{theorem}
\noindent\Kernel{}\quad(\leanname{Jordan.H3OJordan} with the
componentwise helpers \leanname{H3OJordanDiag/X/Y/Z};
\leanname{H2O}, \leanname{H2OProduct}, \leanname{SpinFactor};
\leanname{TraceForm}, \leanname{DVTComplementBridge}.)

\begin{theorem}[DVT two-sided stabilizer action, fully characterized]
\label{thm:dvt}
The DVT/Yokota two-sided action $X \mapsto AXB$ of
$SU(3)\times SU(3)^{\mathrm{op}}$ on the complement of $h_3(\C)$ in
$h_3(\OO)$ has central kernel exactly $\mathbb{Z}_3$ (an iff), the
quotient $(SU(3)\times SU(3)^{\mathrm{op}})/\mathbb{Z}_3$ acts
faithfully, and the image inside the additive automorphisms of the
complement is characterized algebraically: an additive automorphism lies
in the image \emph{iff}, in the $M_3(\C)$ coordinate model, it has the
form $X \mapsto AXB$ with $\det A = \det B = 1$.
\end{theorem}
\noindent\Kernel{}\quad(\leanname{DVTTwoSidedActionKernelZ3Iff};
\leanname{DVTTwoSidedStabilizerMoonshot} for the faithful quotient
action;
\leanname{DVTTwoSidedImageEquiv};
\leanname{DVTFullStabilizerCharacterization} ---
\leanname{dvtQuotientImageSubgroup\_mem\_iff\_twoSidedCompatible}.)

\begin{remark}[Boundary, sharpened against Baez--Schwahn 2026]
This characterizes the \emph{two-sided matrix action} that the DVT route
uses on the complement.  The precise relation to the recent
Baez--Schwahn theorem \cite{BaezSchwahn2026} (verified in primary full
text) is three-part.  \emph{Kernel-checked here:} the coordinate model
of the action of $(SU(3)\times SU(3)^{\mathrm{op}})/\mathbb{Z}_3$ on the
complement of $h_3(\C)$ --- which is, in their notation, the restriction
to the complement of the identity component
$\mathrm{Stab}(B)_0 \cong (SU(3)\times SU(3))/\mathbb{Z}_3$ of the
$h_3(\C)$-stabilizer (their Lemma~4; the identity-component subscript is
load-bearing, since $\mathrm{Stab}(B)$ has a second, antiunitary
component; the $\mathrm{op}$ in our form is notational ---
$SU(3)^{\mathrm{op}} \cong SU(3)$ via $g \mapsto g^{-1}$ --- and records
that the second slot acts by right multiplication in the two-sided
coordinate action $X \mapsto AXB$) --- together with the algebraic
$SU(3)$ of octonion automorphisms preserving $\C\subset\OO$
(Theorem~\ref{thm:su3}).  In their proof shape the two factors play
distinct roles: the first $SU(3)$ acts by conjugation on the ordinary
qutrit $h_3(\C)$ and is the factor their $X$-stabilizer cuts down to
$S(U(2)\times U(1))$, while the second --- the octonionic-automorphism
factor, unimpeded by the $X$-stabilizer --- is the one
Theorem~\ref{thm:su3} realizes; our two kernel-checked theorems are
coordinate avatars of one factor each, not two names for the same
factor.  \emph{External
(source-verified, not formalized):} that this group is all of
$\mathrm{Stab}(B)_0$ inside the compact
$F_4 = \mathrm{Aut}(h_3(\OO))$; the $F_4$-transitivity statements
(their Lemmas~5, 6, 11); and their Theorems~1--2, that for every nested
pair $X \cong h_2(\C) \subset B \cong h_3(\C)$ the intersection
$\mathrm{Stab}(X)\cap\mathrm{Stab}(B)_0$ is $S(U(2)\times U(3))$ ---
their new contribution over Dubois-Violette--Todorov being precisely the
choice-independence via transitivity.  \emph{Open:} a kernel-checked
composition joining these layers; the first missing coordinate rung is
the statement that the $h_2(\C)$-block stabilizer inside the first
$SU(3)$ factor is $S(U(2)\times U(1))$.
\end{remark}

Supporting this, the derivation layer
(\leanname{InnerDerivation},
\leanname{InnerDerivationLieAlgebra},
\leanname{InnerDerivationStandardBSU3Bridge}) verifies that inner
derivations of the coordinate model form a Lie algebra, and bridges the
standard-block subalgebra toward the $\su(3)$ structure used by the DVT
route.

\section{A shared bridge: one complex structure selects the data}
\label{sec:bridge}

The three routes are cohesive because they consume the \emph{same}
choice.  Fixing the imaginary unit $e_{111}$:

\begin{itemize}[leftmargin=2em]
\item it selects the complex line $\C\subset\OO$ and the splitting
  $\OO \cong \C\oplus\C^3$ (Section~\ref{sec:foundation});
\item the automorphisms fixing it are exactly $SU(3)$ acting on the
  $\C^3$ complement (Section~\ref{sec:baez});
\item the ladder operators, idempotent, and minimal ideal of the Furey
  construction are built from the same unit, and their number operators
  produce the charges (Section~\ref{sec:furey});
\item applied entrywise, the splitting embeds $h_3(\C)$ in $h_3(\OO)$
  with the DVT complement action of Section~\ref{sec:jordan};
\item and the gauge endpoint these constructions aim at is the single
  block group $S(U(2)\times U(3))$ with its exact $\mathbb{Z}_6$
  arithmetic (Section~\ref{sec:gauge}).
\end{itemize}

We emphasize --- as the informal literature sometimes does not --- that
these are \emph{parallel uses of a common choice}, not corollaries of one
master theorem.  The formalization makes the parallelism auditable: the
same \leanname{e111}, the same bridge, the same $\C^3$, in every lane.
Whether the parallelism can be upgraded to a dependency chain is the
master question of the next section.

\section{The Jordan--Clifford bridge program: from parallelism to a
master question}
\label{sec:jcbridge}

The honest test of ``unification'' talk in this area is whether one
choice \emph{forces} the rest of the package.  We therefore organize the
frontier as a single master question with graded rungs.  Fix the nested
Jordan flag
\[
X \cong h_2(\C) \;\subset\; B \cong h_3(\C) \;\subset\; J = h_3(\OO).
\]
Does this flag canonically and equivariantly determine, in one dependency
chain: a weak space $W\cong\C^2$ and a color space $V\cong\C^3$; the
gauge group $S(U(W)\times U(V))$; the Jordan complement as
$W$-$V$ bimodule data; the Furey color module as $\Lambda V$; one chiral
generation as $\Lambda^{\mathrm{even}}(W\oplus V)$ (Baez--Huerta's
presentation of the $\mathrm{Spin}(10)$ sixteen \cite{BaezHuerta2009});
and the exact $\mathbb{Z}_6$ kernel of the action on that fermion
module?  Dimension matching is not enough; each arrow must specify its
choices, equivariance, kernel, and uniqueness or residual moduli.

The current rung-by-rung status, with exact trust labels:

\begin{itemize}[leftmargin=2em]
\item \emph{Flag and stabilizer (external glue).}  Baez--Schwahn
  \cite{BaezSchwahn2026} prove (their Theorems~1--2, via transitivity
  Lemmas~5, 6, 11) that for every such flag,
  $\mathrm{Stab}(X)\cap\mathrm{Stab}(B)_0 \cong S(U(2)\times U(3))$ ---
  source-verified with the identity-component subtlety recorded in
  Section~\ref{sec:jordan}; the kernel-checked coordinate avatars of
  both factors are Theorem~\ref{thm:su3} and
  Theorem~\ref{thm:dvt}.
\item \emph{Finite exterior spine (kernel-checked).}  The color
  exterior degrees have dimensions $1+3+3+1=8$ and the even five-mode
  degrees $1+10+5=16$, as genuine \leanname{Module.finrank} theorems
  over Mathlib's exterior powers; the corrected Furey basis is
  \emph{bijective} with three-mode occupancies and all $48$
  creation/annihilation table entries agree exactly with the signed
  exterior rule; and the six-power central kernel is proved equal to
  the injective image of the standard six phases --- with a
  non-membership control --- under the convention $6Y=3N_W-2N_V$,
  which is verified on every five-mode occupation state
  (\leanname{JordanCliffordFermionKernel},
  \leanname{JordanCliffordFureyFockBridge},
  \leanname{JordanCliffordSpinorZ6Bridge}).
\item \emph{Open rungs (the honest gates).}  (i) The whole-submodule
  operator intertwiner: the $48$-entry agreement is a basis-table
  statement; bundling the left multiplications as endomorphisms of the
  ideal and proving a linear-operator intertwining theorem with
  $\Lambda V$ is the decisive next theorem.  (ii) Deriving the
  commuting $W$ and $V$ from the flag rather than positing them.
  (iii) Transporting the finite $\mathbb{Z}_6$ kernel from the
  bidegree arithmetic model to the actual covering-group action on the
  fermion module.  (iv) The chirality question for any
  four-dimensional realization --- an explicit gate, not a footnote.
\end{itemize}

Kill conditions, stated in advance: if the flag supplies only
dimensions, if the weak action must be appended by hand, or if the
$\mathbb{Z}_6$ can be imposed group-theoretically but not recovered
from the fermion action, the correct publication is the precise
noncanonicity theorem, and the unification claim narrows accordingly.
Nothing in this section uses any null-edge result as evidence.

\section{Family symmetry and triality naturality}
\label{sec:family}

Motivated by Furey--Hughes \cite{FureyHughes2025} and nearby $S_3$
family-symmetry models \cite{Gresnigt2026}, the repository provides a
reusable interface rather than a physical claim: if a family action
commutes with the charge and gauge data, it transports eigenvectors with
their eigenvalues, transports charge tables, and leaves every anomaly
coefficient and anomaly-freedom predicate invariant
(\leanname{StandardModel.FamilySymmetry},
\leanname{FamilyAnomalyNaturality} --- the latter proving that
relabeling chiral multiplets changes no anomaly coefficient).  The
interface is instantiated for the triality role permutations of the
Furey--Hughes construction
(\leanname{TrialityFamilyNaturality}: cycle eigenvector transport,
role-constant charge tables, cycle invariance of each charge field, and
the Gell-Mann--Nishijima relation at cycled roles), on top of the
triality role API (\leanname{TrialityTriple},
\leanname{TrialityActionMonoid},
\leanname{TrialityPermutationRepresentation}).

\begin{remark}[Boundary]
This is a conditional naturality package: any concrete triality or
family-symmetry proposal must \emph{prove} the commutation hypotheses
before transporting Standard Model tables or anomaly sums.  It does not
assert that triality is the physical origin of generations, and the full
$\mathrm{tri}(\C)\oplus\mathrm{tri}(\mathbb{H})\oplus\mathrm{tri}(\OO)$
Lie-algebra story of \cite{FureyHughes2025} is not formalized.
\end{remark}

\section{Proof engineering}
\label{sec:eng}

A substantial fraction of the proof scripts was produced with the help of
a Lean proof-search agent (Aristotle, by Harmonic), used as a proposer of
source code, never as an oracle: every accepted proof passed the Lean
kernel in this repository, under the pinned toolchain, and the flagship
surfaces are additionally frozen by the build-enforced guards of
Section~\ref{sec:conventions}.  Agent-produced modules carry provenance
notes, and integration followed a fixed loop: statement prepared and
typechecked locally; proof searched; returned source audited for
semantic alignment (statement unchanged, conventions preserved, no
hidden hypotheses); local kernel check; guard update in the same commit.
Any reader can verify the results by running \leanname{lake build} on the
repository; the kernel, not the agent and not the authors, is the trust
root.

\section{Claim boundaries and non-results}
\label{sec:boundaries}

Verified (all kernel-trust, standard axiom footprint):
finite octonion arithmetic in a fixed convention with a verified bridge;
complex-line and complement structure; the algebraic
$\mathrm{Aut}_{e_{111}}(\OO) \simeq^* SU(3)$ equivalence with Mathlib's
$SU(3)$; Pati--Salam/SM/$SU(5)$ block predicates and their equivalences;
the exact six-element $\mathbb{Z}_6$ kernel, product-level identity-fiber
iff, and unit-level quotient equivalence; Furey ladder $\mathrm{Cl}(6)$
relations, idempotent and ideal arithmetic, charge-operator eigenvalue
table, operator-level Gell-Mann--Nishijima, electroweak $\su(2)$ package,
$Q_{\mathrm{op}}$-derived $U(1)$ anomaly sums, doublet-sector
cancellation, full one-generation anomaly freedom with Witten's
condition, and the one-generation realization package; the Jordan
identity for $h_3(\OO)$, trace-form splitting, and the full two-sided
DVT stabilizer characterization with exact $\mathbb{Z}_3$ kernel;
family-symmetry and triality naturality interfaces.

Not proved, and not claimed:
$\mathrm{Aut}(h_3(\OO)) \cong F_4$;
$\mathrm{Stab}(h_2(\OO)) \cong \mathrm{Spin}(9)$;
the full DVT compact stabilizer-intersection theorem;
any topological, smooth, or connectedness statement about $G_2$,
$SU(3)$, or the gauge quotient;
a derivation of the right-handed fermion sector, of the $W^\pm$
operators, or of the $T_3/Y$ values from the octonionic algebra;
a three-generation representation theorem;
and any statement about dynamics, Lagrangians, masses, mixing, or
Yukawa structure.

\section{Future work}
\label{sec:future}

Near-term targets, in the order we would attempt them:
derive the electroweak $\su(2)$ action from the octonionic operator
algebra (closing the ``defined, not derived'' boundary of
Section~\ref{sec:furey}); upgrade the $\mathbb{Z}_6$ package from the
algebraic to the topological quotient; instantiate the family-symmetry
interface for a concrete Furey--Hughes triality representation beyond
role permutations; and extend the DVT layer from the two-sided action
characterization toward the compact stabilizer statement, which requires
compact-group infrastructure (or a clean workaround) not yet in Mathlib.
A Clifford-envelope layer connecting the Furey operators to Mathlib's
\leanname{CliffordAlgebra} would let the $\mathrm{Cl}(6)$ statements be
stated against library structures rather than concrete matrices.

\appendix
\section{Artifact}
\label{app:artifact}

Toolchain: \leanname{leanprover/lean4:v4.28.0}, Mathlib pinned by the
repository lockfile.  Trusted source relevant to this paper:
\leanname{PhysicsSM/Algebra/Furey} (47 modules, $\approx$11{,}100 lines),
\leanname{PhysicsSM/Algebra/Octonion} (50 modules, $\approx$9{,}100
lines), \leanname{PhysicsSM/Algebra/Jordan} (59 modules,
$\approx$9{,}900 lines), \leanname{PhysicsSM/Gauge} (46 modules,
$\approx$7{,}700 lines), plus the \leanname{StandardModel} family
interface.  Build: \leanname{lake build}; the trusted layer contains no
placeholder tokens and no compiler-trusting evaluation, enforced by
repository-wide scans and, for the flagships
(\leanname{su3Submonoid\_eq\_specialUnitaryGroup},
\leanname{fureyRealizesOneGenerationPackage},
\leanname{alpha1\_nilpotent}, \leanname{anticomm\_1\_1dag}), by
build-failing axiom guards in \leanname{Furey.AxiomGuard}.  The octonion
convention and bridge are cross-validated against an external oracle
script (\leanname{Scripts/oracle/validate\_octonion.py}); the bridge
table is exported in \leanname{Index/convention\_bridge.json}.  A frozen
tagged snapshot with an archival identifier is a release gate to be
completed at submission time.

\begin{thebibliography}{19}

\bibitem{BaezOctonions}
J.~C.~Baez, \emph{The octonions},
Bull.\ Amer.\ Math.\ Soc.\ 39 (2002) 145--205, arXiv:math/0105155.

\bibitem{Baez2021}
J.~C.~Baez, \emph{Can we understand the Standard Model using octonions?},
lecture notes and talks (2021).

\bibitem{BaezHuerta2009}
J.~C.~Baez and J.~Huerta, \emph{The algebra of grand unified theories},
Bull.\ Amer.\ Math.\ Soc.\ 47 (2010) 483--552, arXiv:0904.1556.

\bibitem{Furey2015}
C.~Furey, \emph{Charge quantization from a number operator},
Phys.\ Lett.\ B (2015), arXiv:1603.04078.

\bibitem{Furey2016}
C.~Furey, \emph{Standard model physics from an algebra?},
PhD thesis, University of Waterloo (2015), arXiv:1611.09182.

\bibitem{Furey2018}
C.~Furey, \emph{$SU(3)_C \times SU(2)_L \times U(1)_Y\,(\times U(1)_X)$
as a symmetry of division algebraic ladder operators},
Eur.\ Phys.\ J.\ C 78:375 (2018), arXiv:1806.00612.

\bibitem{FureyHughes2025}
N.~Furey and M.~J.~Hughes, \emph{Three generations and a trio of
trialities}, Phys.\ Lett.\ B 865 (2025) 139473, arXiv:2409.17948.

\bibitem{DVT2018}
M.~Dubois-Violette and I.~Todorov, \emph{Deducing the symmetry of the
standard model from the automorphism and structure groups of the
exceptional Jordan algebra}, Int.\ J.\ Mod.\ Phys.\ A 33 (2018) 1850118,
arXiv:1806.09450.

\bibitem{BaezSchwahn2026}
J.~C.~Baez and P.~Schwahn, \emph{The Standard Model gauge group from
the exceptional Jordan algebra}, arXiv:2606.15235 (2026).

\bibitem{Krasnov2019}
K.~Krasnov, \emph{$SO(9)$ characterisation of the Standard Model gauge
group}, J.\ Math.\ Phys.\ (2021), arXiv:1912.11282.

\bibitem{Krasnov2025}
K.~Krasnov, \emph{Octonions, complex structures, and Standard Model
fermions}, arXiv:2504.16465.

\bibitem{Boyle2020}
L.~Boyle, \emph{The Standard Model, the exceptional Jordan algebra, and
triality}, J.\ Math.\ Phys.\ 67, 071701 (2026), arXiv:2006.16265.

\bibitem{Gresnigt2026}
N.~Gresnigt, \emph{Electroweak structure and three fermion generations
in Clifford algebra with $S_3$ family symmetry}, arXiv:2601.07857.

\bibitem{ToobySmith2025}
J.~Tooby-Smith, \emph{HepLean: digitalising high energy physics},
Comput.\ Phys.\ Commun.\ 308 (2025) 109457, arXiv:2405.08863.

\bibitem{ToobySmith2024}
J.~Tooby-Smith, \emph{Formalization of physics index notation in
Lean~4}, arXiv:2411.07667.

\end{thebibliography}

\end{document}
